\section{Dante’s Empyrean}

In \textit{The Discarded Image}, \textbf{C. S. Lewis} invites the reader to a series of exercises, beginning with looking out at the night sky. When he comes to \textbf{Dante}, however, he realizes that his first exercise was misleading. You do not experience the Empyrean by looking out at the night sky as though it were a canopy encompassing the universe. He writes:

\begin{quotex}
The universe is thus, when our minds are sufficiently freed from the senses, turned inside out.
\end{quotex}


As Fludd's diagram\footnote{\url{https://www.meditationsonthetarot.com/the-cosmic-hierarchy}} indicates, the highest level of the Empyrean is Mens or Consciousness. The world, then, exists in consciousness. In The Metaphysics of Dante's Comedy, \textbf{Christian Moevs} makes a good attempt to describe Dante's worldview in its own terms. Moevs makes these points:

\begin{quotex}
\begin{enumerate}
\item the world of space and time does not itself exist in space and time: it exists in Intellect (the Empyrean, pure conscious being)

\item matter, in medieval hylomorphism, is not something ``material'': it is a principle of unintelligibility, of alienation from conscious being

\item all finite form, that is, all creation, is a self-qualification of Intellect or Being, and only exists insofar as it participates in it

\item Creator and creation are not two, since the latter has no existence independent of the former; but of course creator and creation are not the same

\item God, as the ultimate subject of all experience, cannot be an object of experience: to know God is to know oneself as God, or (if the expression seems troubling) as one ``with'' God or ``in'' God.
\end{enumerate}
\end{quotex}

Obviously, this point of view can also be found in the non-dual teachings of the Eastern Traditions. Of course, this totally alters the meaning of salvation in Dante, or better said, he has an understanding of liberation beyond mere salvation. This is the process of theosis, or becoming God-like. This seems to make Mr. Lewis quite uncomfortable. He writes:

\begin{quotex}
The picture is nothing if not religious. But is the religion in question precisely Christianity?
\end{quotex}

As long as Christianity is assumed to consist of a set of creedal statements to be believed, Mr. Lewis' question is valid. However, if we are describing a conversion, an intellectual conversion\footnote{\url{https://www.gornahoor.net/?p=262}}, the answer is positive: Dante is describing a true picture of the universe. However, as the optical illusion demonstrates, the picture can be seen in two, mutually incompatible, ways. As long as a man is stuck with a worldview of the old lady, he cannot discern the young women. No intellectual argument is persuasive, there needs to be a sudden ``seeing'' of the other view. With that, there is no longer any doubt. But before that conversion, there can be nothing but the demon of dialectics\footnote{\url{https://www.gornahoor.net/?p=813}}. Mr. Moevs explains it this way:

\begin{quotationx}
The kind of enlightenment or revelation this self-professedly salvific text is meant to trigger cannot be reduced to ``doctrine'', to a creed or set of ideas (already familiar to most in any case). To profess a religious or philosophical creed is not to achieve salvation; else many residents of Dante's Hell would live in Paradise, and some who live in Paradise, such as Ripheus, would live in Hell. In the Comedy salvation is rather a self-awakening of the Real to itself in us, the surrender or sacrifice of what we take ourselves and the world to be, a changed experience that is one with a moral transformation. We cannot know what we are until we surrender what we think we are, with all its attendant desires.

By Dante's own definition of philosophy (the reflexive love of intellect for itself), a philosopher whose contribution can be reduced to doctrine, to a new set of ideas, has either been misunderstood or is not very good.
\end{quotationx}


Do we see here yet again the need for the First Trial\footnote{\url{https://www.gornahoor.net/?p=1902}}? We need to sacrifice our personal opinions, the false ideas of who we are, in order to know who we truly are, our True Will. This first requires the Moral Purification of the Will\footnote{\url{https://gornahoor.net/?p=17007}}.

We should not neglect Dante's point of philosophers, of whom he holds \textbf{Aquinas} and \textbf{Aristotle} in high regard. A Cliff Notes summary of their philosophical doctrines is insufficient. What is required is a wrestling with their ideas in order to achieve some sort of intellectual transformation.

The Comedy is all-embracing. As we have been trying to show with Solovyov's Law of Development\footnote{\url{https://www.gornahoor.net/?p=3504}}, the path to self-knowledge requires knowing the world in its Totality, as the total experience of the human race. To pick and choose, to reject a particular period of the cosmic cycle, is to fall short of the task. Moevs describes Dante's view:

\begin{quotex}
In the Comedy, to journey toward self-knowledge is to assimilate as oneself, through direct poetic experience, the entire breadth of human experience in history, in its concrete and particular reality. Hence the Comedy is all encompassing; it is built on human encounters; it is one man's experience, the assimilation of everyone's experience. Through this identification with the totality of possible experience, the obsessive point of view (veduta) and attachments of the individual ego begin to dissolve, and the subject of all experience, reflected as the individual mind, begins to awaken to itself, ultimately to discover itself, on the threshold of the Empyrean, as an extensionless metaphysical point, the light that spawns and is all things. Of that point, nothing can be said: it is not a thing, there is nothing to know. One can only be it; that is to know it.
\end{quotex}



\flrightit{Posted on 2012-01-02 by Cologero}
